\documentclass{article}
\usepackage{amsmath, amssymb, amsthm}
\usepackage{geometry}
\geometry{a4paper, margin=1in}

\title{Answers to Questions on \( C^* \)-Algebras and Extensions}
\author{}
\date{}

\begin{document}

\maketitle

\section*{Answer to: If a linear function is defined and is positive on a self adjoint subspace of a \( C^* \)-algebra, what are the conditions it needs to be satisfied so that it can be extended to a state defined on the whole \( C^* \)-algebra?}

For a positive linear functional \( \phi \) defined on a self-adjoint subspace \( S \) of a \( C^* \)-algebra \( A \) to be extendable to a state on \( A \), the following conditions must be satisfied:

1. \textbf{Positivity}: \( \phi(a) \geq 0 \) for all \( a \in S \) with \( a \geq 0 \).

2. \textbf{Boundedness}: \( \phi \) must be bounded on \( S \), i.e., there exists \( C > 0 \) such that \( |\phi(a)| \leq C \|a\| \) for all \( a \in S \).

3. \textbf{Norm Condition}: If \( S \) contains the identity \( 1 \), then \( \phi(1) = 1 \). If \( S \) does not contain \( 1 \), then \( \phi \) must satisfy \( \phi(a) \leq \|a\| \) for all \( a \in S \).

These conditions ensure that \( \phi \) can be extended to a state \( \tilde{\phi} \) on \( A \).

---

\section*{Answer to: Doesn't the linear combination of \( p \) and \( q \) belong to the subspace by definition of a subspace, thus the subspace contains the identity, which contradicts the original assumption?}

You are correct! If \( p \) and \( q \) are projectors in \( S \) such that \( p + q = I \), then \( S \) contains the identity \( I \). This contradicts the assumption that \( S \) does not contain the identity. Therefore, the earlier example was flawed. Let me provide a corrected analysis.

---

\section*{Answer to: Is \( \phi(p) \leq 1 \) for all projectors \( p \in S \), when \( S \) does not contain the identity, the necessary and sufficient condition for \( \phi \) to be extended to a state on all of \( A \)?}

The condition \( \phi(p) \leq 1 \) for all projectors \( p \in S \) is \textbf{necessary} but \textbf{not sufficient} for \( \phi \) to be extended to a state on \( A \). Here’s why:

1. \textbf{Necessary Condition}: If \( \phi \) can be extended to a state \( \tilde{\phi} \) on \( A \), then \( \phi(p) = \tilde{\phi}(p) \leq 1 \) for all projectors \( p \in S \).

2. \textbf{Not Sufficient}: The condition \( \phi(p) \leq 1 \) does not guarantee that \( \phi \) is positive or bounded on all of \( S \). For example, \( \phi \) could grow arbitrarily large on other elements of \( S \), violating the boundedness condition.

Thus, while \( \phi(p) \leq 1 \) is necessary, it is not sufficient for extension to a state.

---

\section*{Answer to: Is the set of positive elements \( a \in S \) a convex set?}

Yes, the set of positive elements \( S_+ = \{ a \in S \mid a \geq 0 \} \) is a \textbf{convex set}. This is because:

1. If \( a, b \in S_+ \), then \( \lambda a + (1 - \lambda) b \in S_+ \) for all \( \lambda \in [0, 1] \).

2. The convex combination of positive elements is positive, as the set of positive elements in a \( C^* \)-algebra is a convex cone.

---

\section*{Answer to: What are the extreme points of the convex set of positive elements of \( S \)?}

The \textbf{extreme points} of the convex set \( S_+ \) are the positive elements in \( S \) that cannot be written as non-trivial convex combinations of other positive elements in \( S \). These extreme points are closely related to \textbf{minimal projections} in \( S \):

1. If \( S \) contains minimal projections, they are extreme points of \( S_+ \).

2. If \( S \) does not contain minimal projections, the extreme points may correspond to other indecomposable positive elements.

---

\section*{Answer to: If the \( C^* \)-algebra is a W\(^*\)-algebra, does this change anything?}

If \( A \) is a W\(^*\)-algebra (von Neumann algebra), the analysis of \( S_+ \) and its extreme points is influenced by the additional structure of the W\(^*\)-algebra:

1. The weak operator topology allows us to study limits of nets of positive elements.

2. Minimal projections in \( S \) (if they exist) are extreme points of \( S_+ \).

3. The type of the W\(^*\)-algebra (Type I, II, or III) affects the structure of extreme points.

---

\section*{Answer to: If \( S \) is a self-adjoint subspace of a W\(^*\)-algebra, how does this affect the preceding answers?}

If \( S \) is a self-adjoint subspace of a W\(^*\)-algebra \( A \), the analysis of \( S_+ \) and its extreme points is influenced by the W\(^*\)-algebra structure:

1. Minimal projections in \( S \) (if they exist) are extreme points of \( S_+ \).

2. The weak operator topology and the type of the W\(^*\)-algebra affect the structure of extreme points.

3. If \( S \) contains the identity, the extreme points of the set of normal states on \( S \) are related to the extreme points of \( S_+ \).

---

\section*{Answer to: Do the inf and sup values of \( \phi \) on \( S_+ \) correspond to extreme points of \( S_+ \)?}

The \textbf{supremum} of \( \phi \) on \( S_+ \) is often achieved at an extreme point of \( S_+ \), especially if \( S_+ \) is compact. However, the \textbf{infimum} of \( \phi \) on \( S_+ \) is less directly tied to extreme points:

1. The supremum is often achieved at an extreme point.

2. The infimum may or may not be achieved at an extreme point, particularly if \( \phi \) is positive (in which case the infimum is \( 0 \)).

---

\section*{Answer to: If \( S_+ \) in a W\(^*\)-algebra is weakly compact, are its extreme points minimal projectors?}

If \( S_+ \) is weakly compact in a W\(^*\)-algebra:

1. \textbf{Minimal projections} in \( S \) are always \textbf{extreme points} of \( S_+ \).

2. Whether \textbf{all extreme points} of \( S_+ \) are minimal projections depends on the structure of \( S \):
   - If \( S \) contains all minimal projections or is a face of the positive cone, then the extreme points of \( S_+ \) are precisely the minimal projections.
   - Otherwise, the extreme points may include other indecomposable positive elements.

\end{document}